\hspace{5cm} \textbf{\Huge Résumé}

\vspace{1cm}

Version longue : 

La théorie de l’Inflation constitue aujourd’hui le socle du Modèle Standard de la cosmologie. Toutefois, aucune preuve observationnelle directe n’en a encore été apportée, et démontrer l’existence de cette phase d’expansion demeure l’un des défis majeurs de la cosmologie observationnelle moderne. Les expériences actuelles et futures visent à détecter les modes B primordiaux de polarisation du Fond Diffus Cosmologique, prédits par l’Inflation, afin d’en établir l’existence et d’en contraindre les paramètres.

L’une de ces expériences, QUBIC, relève ce défi en adoptant une architecture instrumentale originale fondée sur l’interférométrie bolométrique. Cette approche permet un contrôle particulièrement fin des effets systématiques grâce à la technique de self-calibration, tout en offrant une amélioration de la résolution fréquentielle via l’imagerie spectrale. Cette dernière propriété est essentielle pour séparer le signal primordial des différentes sources de contamination, notamment les avant-plans astrophysiques et les émissions atmosphérique.

Les travaux présentés dans cette thèse portent principalement sur le développement du code de simulation instrumentale, ainsi que sur la conception d’algorithmes de reconstruction de cartes et de séparation de composantes astrophysiques. Je présente d’abord mes contributions à la source de calibration de QUBIC et à son système GPS associé, en vue de préparer les opérations de self-calibration. Je détaille ensuite les développements réalisés sur le code de simulation de l’instrument.

La majeure partie de cette thèse est consacrée aux algorithmes de reconstruction de cartes et de séparation de composantes. J’y décris en particulier mon implication dans le développement des méthodes de Frequency Map-Making et de Component Map-Making. Je présente également les améliorations apportées au pipeline de simulation afin d’en accroître le réalisme, ainsi que le développement d’une approche innovante fondée sur les réseaux de neurones pour résoudre le problème inverse du map-making. Ces différentes méthodes sont ensuite appliquées à l’atténuation de la contamination atmosphérique.

Enfin, je conclus en présentant des prévisions comparatives des différents algorithmes, sous diverses hypothèses d’avant-plans astrophysiques et de configurations instrumentales, afin d’évaluer la sensibilité de QUBIC à la détection des modes B primordiaux.

Version courte : 

La théorie de l’Inflation constitue le socle du Modèle Standard de la cosmologie. Toutefois, aucune preuve observationnelle directe n’a encore été établie, et la détection des modes B primordiaux de polarisation du Fond Diffus Cosmologique demeure un enjeu central pour en confirmer l’existence et en contraindre les paramètres.

L’expérience QUBIC propose une approche originale fondée sur l’interférométrie bolométrique, permettant un contrôle fin des effets systématiques grâce à la self-calibration et une meilleure résolution fréquentielle via l’imagerie spectrale. Cette dernière est essentielle pour séparer le signal primordial des contaminations astrophysiques et atmosphériques.

Cette thèse porte principalement sur le développement du code de simulation instrumentale et sur la conception d’algorithmes de reconstruction de cartes et de séparation de composantes. Après une contribution à la source de calibration et à son système GPS, je détaille les améliorations apportées à la simulation. La majeure partie du travail concerne le développement des méthodes de Frequency Map-Making et de Component Map-Making, l’amélioration du réalisme des simulations, ainsi qu’une approche innovante du problème inverse du map-making fondée sur les réseaux de neurones. Ces outils sont appliqués à l’atténuation de la contamination atmosphérique.

Enfin, je présente des prévisions comparatives, sous différentes hypothèses d’avant-plans et de configurations instrumentales, afin d’évaluer la sensibilité de QUBIC à la détection des modes B primordiaux.

\vspace{0.5cm}

\textbf{Mots clés :} Fond diffus cosmologique - Inflation - Interférometrie Bolométrique - Imagerie Spectrale - Séparation de composantes - Map-Making

\newpage

\hspace{5cm} \textbf{\Huge Abstract}

\vspace{1cm}

\vspace{0.5cm}

\textbf{Keywords :} Cosmic Microwave Background - Inflation - Bolometric Interferometry - Spectral-Imaging - Component Separation - Map-Making
